\subsection{Affine Independence and Barycentric Coordinates}

Throughout this subsection, $\mathcal{V}$ is finite-dimensional over a fixed field $\mathbb{F}$.

\begin{definition} \label{def:affinely-independent-set} \hlt{Affinely Independent Sets}\\
A set $E\subseteq\mathcal{V}$ is \hlt{affinely independent} iff $E\neq\emptyset$ and $(\forall\,\mathbf{x}\in E)\,(E-\mathbf{x})\setminus\{\mathbf{0}\}$ is linearly independent.
\end{definition}

\begin{theorem} \label{thm:affine-independence-equivalences} \hlt{Affine Independence Criteria}\\
Let $\emptyset\neq E\subseteq\mathcal{V}$ and, for $F\subseteq\mathcal{V}$ and $n\in\mathbb{Z}_{>0}$, put $F_{\neq}^n\defeq\{\mathbf{x}\in F^n:\mathbf{x}_j\neq\mathbf{x}_k\ (1\leq j<k\leq n)\}$.\\
The following conditions are equivalent.
\begin{enumerate}[label=(\roman*)]
\setlength{\itemsep}{0pt}
\item $E$ is affinely independent.
\item $(\forall\,\mathbf{x}\in E)\,\mathbf{x}\notin\operatorname{affhull}(E\setminus\{\mathbf{x}\})$, where $\operatorname{affhull}(\emptyset)=\emptyset$.
\item $(\forall\,n\in\mathbb{Z}_{>0})(\forall\,(\mathbf{x}_1,\ldots,\mathbf{x}_n)\in E_{\neq}^n)(\forall\,\mathbf{r}\in\mathbb{F}^n)\,[\sum\limits_{j=1}^n r_j\mathbf{x}_j=\mathbf{0}\land\sum\limits_{j=1}^n r_j=0\Rightarrow\mathbf{r}=\mathbf{0}]$.
\item $(\forall\,n\in\mathbb{Z}_{>0})(\forall\,(\mathbf{x}_1,\ldots,\mathbf{x}_n)\in E_{\neq}^n)(\forall\,\mathbf{r},\mathbf{s}\in\mathbb{F}^n)\,[\sum\limits_{j=1}^n r_j\mathbf{x}_j=\sum\limits_{j=1}^n s_j\mathbf{x}_j\land\sum\limits_{j=1}^n r_j=\sum\limits_{j=1}^n s_j=1\Rightarrow\mathbf{r}=\mathbf{s}]$.
\end{enumerate}
\begin{proof}
\begin{enumerate}[label=(\roman*)]
\setlength{\itemsep}{0pt}
\item[$(i)$]$\Rightarrow (ii)$ Fix $\mathbf{x}\in E$. If $E=\{\mathbf{x}\}$, then $\mathbf{x}\notin\operatorname{affhull}(\emptyset)$ by the empty-hull convention. Suppose $E\neq\{\mathbf{x}\}$ and $\mathbf{x}\in\operatorname{affhull}(E\setminus\{\mathbf{x}\})$.\\
By Theorem \ref{thm:affine-hull-span-formula}(i), $\mathbf{x}\in\operatorname{affspan}(E\setminus\{\mathbf{x}\})$. Collect the repeated points in an affine representation of $\mathbf{x}$.\\
Fix $m\in\mathbb{Z}_{>0}$ and $(\mathbf{y}_1,\ldots,\mathbf{y}_m)\in(E\setminus\{\mathbf{x}\})_{\neq}^m$ as its distinct points.\\
The collected coefficients form $\mathbf{r}\in\mathbb{F}^m$ with $\mathbf{x}=\sum\limits_{j=1}^m r_j\mathbf{y}_j$ and $\sum\limits_{j=1}^m r_j=1$.\\
The relation $\sum\limits_{j=1}^m r_j(\mathbf{y}_j-\mathbf{x})=\mathbf{0}$ with $\mathbf{r}\neq\mathbf{0}$ holds by these identities, contradicting (i). 
\item[$(ii)$]$\Rightarrow (iii)$ Fix $n\in\mathbb{Z}_{>0}$ and $(\mathbf{x}_1,\ldots,\mathbf{x}_n)\in E_{\neq}^n$. Let $\mathbf{r}\in\mathbb{F}^n$ satisfy $\sum\limits_{j=1}^n r_j\mathbf{x}_j=\mathbf{0}$ and $\sum\limits_{j=1}^n r_j=0$.\\
If $\mathbf{r}\neq\mathbf{0}$, choose $k$ with $r_k\neq0$. The inequality $n\geq2$ holds by the contradiction $r_k=\sum\limits_{j=1}^n r_j=0$ when $n=1$.\\
The identities $\mathbf{x}_k=\sum\limits_{j\neq k}(-\frac{r_j}{r_k})\mathbf{x}_j$ and $\sum\limits_{j\neq k}-\frac{r_j}{r_k}=1$ hold by rearrangement.\\
$\mathbf{x}_k\in\operatorname{affhull}(E\setminus\{\mathbf{x}_k\})$ by Theorem \ref{thm:affine-hull-span-formula}(i), contrary to (ii). $\mathbf{r}=\mathbf{0}$ by contradiction. 
\item[$(iii)$]$\Rightarrow (iv)$ Subtraction produces $\sum\limits_{j=1}^n(r_j-s_j)\mathbf{x}_j=\mathbf{0}$ and $\sum\limits_{j=1}^n(r_j-s_j)=0$. Thus $\mathbf{r}=\mathbf{s}$ by (iii). \item[$(iv)$]$\Rightarrow (i)$ Fix $\mathbf{x}\in E$, $m\in\mathbb{Z}_{>0}$, and $(\mathbf{y}_1,\ldots,\mathbf{y}_m)\in(E\setminus\{\mathbf{x}\})_{\neq}^m$. Let $\mathbf{a}\in\mathbb{F}^m$ satisfy $\sum\limits_{j=1}^m a_j(\mathbf{y}_j-\mathbf{x})=\mathbf{0}$.\\
The tuple $(\mathbf{y}_1,\ldots,\mathbf{y}_m,\mathbf{x})$ lies in $E_{\neq}^{m+1}$. Put $\mathbf{c}\defeq(a_1,\ldots,a_m,-\sum\limits_{j=1}^m a_j)$ and $\mathbf{b}\defeq(0,\ldots,0,1)$.\\
Both $\mathbf{b}$ and $\mathbf{b}+\mathbf{c}$ have coefficient sum $1$ and represent $\mathbf{x}$, so $\mathbf{c}=\mathbf{0}$ by (iv).\\
$\mathbf{a}=\mathbf{0}$ by the definition of $\mathbf{c}$, and $E$ is affinely independent by the arbitrariness of all choices.
\end{enumerate}
\hlt{Structural Commentary}: Pass from excluded-point hulls to zero-sum relations, then identify them with coordinate differences.\\
\hlt{Key Insight}: Affine representations differ by a zero-sum linear relation.
\end{proof}
\end{theorem}

\begin{corollary} \label{cor:finite-affine-independence-dimension-criterion} \hlt{Dimension Criterion for Affine Independence}\\
Let $E\subseteq\mathcal{V}$ be finite and nonempty with $\abs{E}=n$. Then $E$ is affinely independent iff $\dim(\operatorname{affhull}(E))=n-1$.
\begin{proof}
\begin{enumerate}[label=\roman*.]
\setlength{\itemsep}{0pt}
\item[$(\Rightarrow)$] Fix $\mathbf{x}\in E$ and put $B_{\mathbf{x}}\defeq(E-\mathbf{x})\setminus\{\mathbf{0}\}$.\\
The set $B_{\mathbf{x}}$ has $n-1$ elements, is linearly independent by the forward hypothesis, and spans $\operatorname{span}(E-\mathbf{x})$.\\
The identity $\dim(\operatorname{affhull}(E))=\dim(\operatorname{span}(E-\mathbf{x}))=n-1$ holds by Theorem \ref{thm:affine-hull-span-formula}(ii).
\item[$(\Leftarrow)$] Fix $\mathbf{x}\in E$ and put $B_{\mathbf{x}}\defeq(E-\mathbf{x})\setminus\{\mathbf{0}\}$. By Theorem \ref{thm:affine-hull-span-formula}(ii), $\dim(\operatorname{span}(E-\mathbf{x}))=n-1=\abs{B_{\mathbf{x}}}$.\\
Any ordering of $B_{\mathbf{x}}$ is a basis of $\operatorname{span}(E-\mathbf{x})$ by Proposition \ref{prop:full-dimension-basis-criteria}(i).\\
$B_{\mathbf{x}}$ is linearly independent by the basis property, and $E$ is affinely independent by the arbitrariness of $\mathbf{x}$.
\end{enumerate}
\hlt{Structural Commentary}: Compare each translated difference set's cardinality with the dimension of its span.\\
\hlt{Key Insight}: Translation removes one point but preserves all affine directions.
\end{proof}
\end{corollary}

\begin{definition} \label{def:affine-basis-barycentric-coordinates} \hlt{Affine Bases and Barycentric Coordinates}\\
Let $A=\mathbf{x}+\mathcal{S}$ be an $n$-flat.
\begin{enumerate}[label=(\roman*)]
\setlength{\itemsep}{0pt}
\item A finite subset $\mathcal{B}\subseteq A$ is an \hlt{affine basis} of $A$ iff $\mathcal{B}$ is affinely independent and $\operatorname{affspan}(\mathcal{B})=A$.
\item A tuple $(\mathbf{p}_1,\ldots,\mathbf{p}_m)$ is an \hlt{ordered affine basis} of $A$ iff its entries are distinct and $\{\mathbf{p}_1,\ldots,\mathbf{p}_m\}$ is an affine basis of $A$.
\item Let $(\mathbf{p}_1,\ldots,\mathbf{p}_m)$ be an ordered affine basis of $A$ and $\mathbf{v}\in A$.\\
A vector $\mathbf{r}=(r_1,\ldots,r_m)\in\mathbb{F}^m$ is a \hlt{barycentric coordinate vector} of $\mathbf{v}$ iff $\mathbf{v}=\sum\limits_{j=1}^{m}r_j\mathbf{p}_j$ and $\sum\limits_{j=1}^{m}r_j=1$.
\end{enumerate}
\end{definition}

\begin{theorem} \label{thm:affine-bases-and-barycentric-coordinate-uniqueness} \hlt{Affine Bases and Barycentric Coordinates}\\
Let $A=\mathbf{x}+\mathcal{S}$ be an $n$-flat.
\begin{enumerate}[label=(\roman*)]
\setlength{\itemsep}{0pt}
\item If $(\mathbf{s}_1,\ldots,\mathbf{s}_n)$ is an ordered basis of $\mathcal{S}$, then $(\mathbf{x}+\mathbf{s}_1,\ldots,\mathbf{x}+\mathbf{s}_n,\mathbf{x})$ is an ordered affine basis of $A$.
\item If $\mathcal{B}$ is an affine basis of $A$, then $\abs{\mathcal{B}}=n+1$.
\item If $(\mathbf{p}_1,\ldots,\mathbf{p}_{n+1})$ is an ordered affine basis of $A$, then $(\forall\,\mathbf{v}\in A)(\exists!\,\mathbf{r}\in\mathbb{F}^{n+1})\,[\mathbf{v}=\sum\limits_{j=1}^{n+1}r_j\mathbf{p}_j\land\sum\limits_{j=1}^{n+1}r_j=1]$.
\end{enumerate}
\begin{proof}
\begin{enumerate}[label=(\roman*)]
\setlength{\itemsep}{0pt}
\item Put $\mathbf{p}_j\defeq\mathbf{x}+\mathbf{s}_j$ for $1\leq j\leq n$ and $\mathbf{p}_{n+1}\defeq\mathbf{x}$. The entries of $(\mathbf{p}_1,\ldots,\mathbf{p}_{n+1})$ are distinct by linear independence of $(\mathbf{s}_1,\ldots,\mathbf{s}_n)$.\\
Every zero-sum relation on a subtuple extends by zero coefficients to the full tuple.\\
For $\mathbf{a}\in\mathbb{F}^{n+1}$ with $\sum\limits_{j=1}^{n+1}a_j\mathbf{p}_j=\mathbf{0}$ and $\sum\limits_{j=1}^{n+1}a_j=0$, the relation $\sum\limits_{j=1}^{n}a_j\mathbf{s}_j=\mathbf{0}$ holds by substitution.\\
The identities $a_1=\cdots=a_n=0$ and $a_{n+1}=0$ hold by independence of the $\mathbf{s}_j$ and the coefficient sum, respectively.\\
The set $\{\mathbf{p}_1,\ldots,\mathbf{p}_{n+1}\}$ is affinely independent by Theorem \ref{thm:affine-independence-equivalences}(iii).\\
By Theorem \ref{thm:affine-hull-span-formula}(i), $\operatorname{affspan}(\{\mathbf{p}_1,\ldots,\mathbf{p}_{n+1}\})=\mathbf{x}+\operatorname{span}(\mathbf{s}_1,\ldots,\mathbf{s}_n)=A$.\\
The tuple $(\mathbf{p}_1,\ldots,\mathbf{p}_{n+1})$ is an ordered affine basis of $A$ by the preceding distinctness, independence, and spanning identities.
\item Put $m\defeq\abs{\mathcal{B}}$. By Corollary \ref{cor:finite-affine-independence-dimension-criterion}, $\dim(\operatorname{affhull}(\mathcal{B}))=m-1$. The identity $\operatorname{affhull}(\mathcal{B})=A$ holds by $\operatorname{affspan}(\mathcal{B})=A$ and Theorem \ref{thm:affine-hull-span-formula}(i).\\
The equality $\abs{\mathcal{B}}=m=n+1$ holds by $n=\dim(A)=\dim(\operatorname{affhull}(\mathcal{B}))=m-1$.
\item Fix $\mathbf{v}\in A$ and put $\mathcal{P}\defeq\{\mathbf{p}_1,\ldots,\mathbf{p}_{n+1}\}$. $\mathbf{v}\in\operatorname{affspan}(\mathcal{P})$ by $\operatorname{affspan}(\mathcal{P})=A$.\\
Choose $\ell\in\mathbb{Z}_{>0}$, $\mathbf{q}\in\mathcal{P}^{\ell}$, and $\mathbf{c}\in\mathbb{F}^{\ell}$ from this affine representation.\\
The chosen data satisfy $\mathbf{v}=\sum\limits_{k=1}^{\ell}c_k\mathbf{q}_k$ and $\sum\limits_{k=1}^{\ell}c_k=1$. For $1\leq j\leq n+1$, put $I_j\defeq\{k\in\{1,\ldots,\ell\}:\mathbf{q}_k=\mathbf{p}_j\}$ and $r_j\defeq\sum\limits_{k\in I_j}c_k$.\\
The sets $I_1,\ldots,I_{n+1}$ partition $\{1,\ldots,\ell\}$. The identities $\mathbf{v}=\sum\limits_{j=1}^{n+1}r_j\mathbf{p}_j$ and $\sum\limits_{j=1}^{n+1}r_j=1$ hold by this partition.\\
The vector $\mathbf{r}\defeq(r_1,\ldots,r_{n+1})$ is unique by Theorem \ref{thm:affine-independence-equivalences}(iv).
\end{enumerate}
\hlt{Structural Commentary}: Adjoin an affine origin, compare dimensions, then group coefficients by basis point.\\
\hlt{Key Insight}: Barycentric uniqueness is the affine-independence uniqueness criterion.
\end{proof}
\end{theorem}

\begin{example} \label{ex:barycentric-coordinates-on-plane} \hlt{Barycentric Coordinates on a Plane}\\
Let $\mathbf{x},\mathbf{v}_1,\mathbf{v}_2\in\mathbb{R}^3$ with $(\mathbf{v}_1,\mathbf{v}_2)$ linearly independent, and put $A\defeq\mathbf{x}+\operatorname{span}(\mathbf{v}_1,\mathbf{v}_2)$. Put $\mathbf{p}_1\defeq\mathbf{x}+\mathbf{v}_1$, $\mathbf{p}_2\defeq\mathbf{x}+\mathbf{v}_2$, and $\mathbf{p}_3\defeq\mathbf{x}$.\\
The tuple $(\mathbf{p}_1,\mathbf{p}_2,\mathbf{p}_3)$ is an ordered affine basis of $A$ by Theorem \ref{thm:affine-bases-and-barycentric-coordinate-uniqueness}(i). For $(t_1,t_2)\in\mathbb{R}^2$, put $\mathbf{a}\defeq\mathbf{x}+t_1\mathbf{v}_1+t_2\mathbf{v}_2$.\\
Then $\mathbf{a}=t_1\mathbf{p}_1+t_2\mathbf{p}_2+(1-t_1-t_2)\mathbf{p}_3$ and $t_1+t_2+(1-t_1-t_2)=1$.\\
Thus $(t_1,t_2,1-t_1-t_2)$ is the barycentric coordinate vector of $\mathbf{a}$ by Theorem \ref{thm:affine-bases-and-barycentric-coordinate-uniqueness}(iii).
\end{example}
